\section{Evaluation Rules and Restrictions}

\subsection*{Truthiness}
scm-slang follows Scheme truthiness: only \texttt{\#f} is false. All other
values are true.

\subsection*{Definitions and Mutation}
\begin{itemize}
\item \texttt{define} introduces a new binding in the current environment.
\item \texttt{set!} requires that the identifier already be defined; otherwise
      it is an error.
\item Redefinition is permitted if the identifier is already in scope; scm-slang
      does not block shadowing across nested scopes created by \texttt{lambda} or
      \texttt{let}.
\end{itemize}

\subsection*{Sequencing}
\texttt{begin} evaluates its subexpressions in order and returns the value of the
last expression. A \texttt{begin} form does not introduce a new scope.

\subsection*{Conditionals}
\texttt{if} accepts two or three operands. If the alternative is omitted, the
value is \texttt{undefined}. \texttt{cond} evaluates predicates in order and
returns the first consequent whose predicate is true. An \texttt{else} clause, if
present, must appear last.

\subsection*{Delayed Evaluation}
\texttt{delay} produces a promise. \texttt{force} evaluates a promise at most once
and memoizes the result.

\subsection*{Quotation}
Quoted data are not evaluated. Quasiquotation supports unquote and
unquote-splicing; unquote and splicing are only valid inside a quasiquoted
context.

\subsection*{Imports and Exports}
\texttt{import} and \texttt{export} follow a scm-slang specific syntax that is
compatible with the Source Academy module system. Imports must appear at the
top level.
